% =============================================================================
% Clifford Seminar Group — Sitzungsprotokoll Template
%
% How to use this template:
%   1. Copy this file to:
%        tex/sitzungsprotokoll/sessionNN-<slug>/sessionNN-sitzungsprotokoll-YYYY-MM-DD.tex
%   2. Fill in the session metadata card (\seminartitle, \seminartopicref,
%      \seminarreferat, \seminarscribe, \seminarpresent, etc.).
%   3. Compile with:
%        latexmk -pdf sessionNN-sitzungsprotokoll-YYYY-MM-DD.tex
%      (If the protocol contains no citations, a single pdflatex pass twice
%       is sufficient; pass [sitzungsprotokoll,nobib] if biblatex is unused.)
%
% Relative paths to clifford-seminar.sty, assets/logo.jpg, and
% bibliography/clifford-seminar-group.bib work identically here and in any
% tex/sitzungsprotokoll/sessionNN-<slug>/ directory.
% =============================================================================
\documentclass[11pt,a4paper]{article}

\usepackage[sitzungsprotokoll]{../../templates/clifford-seminar}

% -----------------------------------------------------------------------------
% Session Metadata (Edit for each meeting)
% -----------------------------------------------------------------------------
\seminartitle[Session NN Record]{%
  Universal Clifford Construction and the Geometric Product%
}
\seminarsubtitle{Template Session Record --- Replace with Session Focus}
\seminarsession{Session NN}
\seminarmodule{Module II}
\seminardate{DD Month YYYY}
\seminarinstitution{Universidad Distrital Francisco Jos\'e de Caldas}

\seminartopicref{Thesenpapier NN --- Universal Clifford Construction, Involutions, and the Geometric Product}
\seminarreferat{Senior Name \& Junior Name}
\seminarkorreferat{Respondent Name}
\seminarscribe{Scribe Name}
\seminarpresent{Participant A, Participant B, Participant C, Participant D}
\seminaronline{Participant E, Participant F} % Leave empty if no online attendees

\begin{document}

\makeseminarheader

% =============================================================================
\section{Recap and Opening Context}
\label{sec:recap}
% =============================================================================

Begin with a brief synthesis of how the meeting opened: continuity from the
previous session, any unresolved exercise revisited at the blackboard, and
the core question framing the day's derivations.

We opened by recalling from Session~01 that a symmetric bilinear form $b$ on
a real linear space $V$ induces the musical isomorphism $b^\flat : V \to V^*$,
and that the free tensor algebra $T(V)$ imposes no metric relations among its
generators. The guiding question for this session was how to quotient $T(V)$
by the quadratic form $q(v) = b(v, v)$ without losing associativity.

% =============================================================================
\section{Blackboard Derivations (Referat)}
\label{sec:derivations}
% =============================================================================

Document the live blackboard derivations faithfully, highlighting logical
transitions, key equations, and moments where the group paused to verify
hypotheses or correct signs.

\subsection{Quotient Construction and Anticommutation}
\label{subsec:quotient}

Starting from the ideal $\mathcal{I}_q = ( v \otimes v - q(v) 1 )
\subset T(V)$, the speakers formed the quotient $\Cl(V,q) = T(V)/\mathcal{I}_q$
and applied the relation $(u + v)^2 = q(u + v) 1$ to obtain
\begin{equation}
  \label{eq:proto-polarization}
  uv + vu = 2\, b(u, v) 1.
\end{equation}
Two immediate consequences were noted at the board:
\begin{itemize}
  \item Orthogonal vectors ($b(u,v) = 0$) strictly anticommute: $uv = -vu$.
  \item Any non-null vector ($q(v) \neq 0$) admits a unique multiplicative
    inverse given by $v^{-1} = v / q(v)$.
\end{itemize}

\begin{discussionnote}[Board Clarification --- Invertibility of Null Vectors]
  A participant asked whether $v^{-1} = v / q(v)$ fails in Minkowski
  spacetime $\R^{1,3}$. The group verified that lightlike (null) vectors
  satisfy $q(v) = 0$ and therefore form zero divisors ($v^2 = 0$) rather than
  invertible elements. This distinction will become essential when analyzing
  null bivectors and electromagnetic plane waves in Module~IV.
\end{discussionnote}

\subsection{Decomposition of the Geometric Product}
\label{subsec:decomp}

Splitting the product of two vectors into symmetric and antisymmetric parts
yielded
\[
  uv = \underbrace{\frac{1}{2}(uv + vu)}_{u \cdot v}
     + \underbrace{\frac{1}{2}(uv - vu)}_{u \wedge v}.
\]
In $\Cl(3,0)$ with orthonormal basis $\{e_1, e_2, e_3\}$, the speakers verified
explicitly that $(e_1 e_2)^2 = e_1 e_2 e_1 e_2 = -e_1^2 e_2^2 = -1$, showing
that unit spatial bivectors act algebraically as generators of rotations
without introducing complex scalars.

% =============================================================================
\section{Respondent Critique (Korreferat) and Plenary Debate}
\label{sec:debate}
% =============================================================================

Summarize the respondent's critical observations and the plenary discussion
structured around the theses from the \emph{Thesenpapier}.

\paragraph{Thesis~I (Geometric Product vs.\ Cross Product).}
The respondent compared the bivector $u \wedge v$ with the traditional vector
cross product $\vec{u} \times \vec{v}$ in $\R^3$, noting that the two are
linked only through Hodge duality with the Euclidean pseudoscalar
$I_3 = e_1 e_2 e_3$:
\[
  u \wedge v = I_3\, (u \times v).
\]
Plenary discussion emphasized that while $u \times v$ is restricted to three
dimensions and flips sign spuriously under spatial inversion (axial vs.\ polar
vectors), the bivector $u \wedge v$ is well-defined in any dimension and
transforms naturally under reflections.

\paragraph{Pedagogical Note on Grading.}
There was brief confusion over why $\Cl(V, q)$ is only $\Z_2$-graded as an
algebra while still having blades of grade $0, 1, \dots, n$. It was clarified
that multiplying two grade-$1$ vectors produces a sum of grade-$0$ and
grade-$2$ terms, preserving parity (even/odd) rather than strict integer
degree \cite[Sect.~2.2]{vazIntroductionCliffordAlgebras2016}.

% =============================================================================
\section{Conclusions, Open Problems, and Next Session}
\label{sec:conclusions}
% =============================================================================

\begin{takeaways}[Resolved Conclusions]
  \begin{enumerate}
    \item The Clifford relation $v^2 = q(v) 1$ uniquely couples the metric
      structure of $(V, q)$ to an associative algebra product via
      \cref{eq:proto-polarization}.
    \item Non-null vectors are invertible in $\Cl(V, q)$ via $v^{-1} = v/q(v)$,
      whereas null vectors in indefinite signatures are zero divisors.
    \item Oriented planes are represented directly by bivectors $u \wedge v$,
      avoiding dimension-specific cross products.
  \end{enumerate}
\end{takeaways}

\begin{nextsession}[Open Questions \& Preparation for Next Session]
  \begin{itemize}
    \item \textbf{Unfinished item:} Complete the universal-property proof of
      the grade involution $\ginv{u}$ and reversion $\rev{u}$ (deferred due to
      time).
    \item \textbf{Exercises to prepare:} Problem~2(b) ($\Cl^+(3,0) \cong
      \mathbb{H}$) and Problem~3 from Thesenpapier~NN.
    \item \textbf{Next session responsibility:} Assign the next protocol scribe
      and confirm preparatory meeting date with the seminar leader.
  \end{itemize}
\end{nextsession}

% Uncomment the following line if this protocol cites works from the shared .bib file:
\printbibliography[title={References Cited in Discussion}]

\end{document}
